
\section*{Глава 2. Проектирование веб-приложения}
\addcontentsline{toc}{section}{Глава 2. Проектирование веб-приложения}
\stepcounter{section}


\subsection{Анализ целевой аудитории}

Основными пользователями разрабатываемой системы являются покупатели спортивных
товаров и администраторы интернет-магазина. Выделяются две ключевые роли:

\begin{itemize}
    \item \textbf{Покупатель} --- зарегистрированный пользователь, осуществляющий
          поиск и выбор товаров, оформление заказов, а также управление личным
          кабинетом и просмотр истории покупок.
    \item \textbf{Администратор} --- сотрудник магазина, ответственный за
          управление каталогом товаров, обработку и отслеживание статусов
          заказов, а также работу с учётными записями пользователей.
\end{itemize}

Типичный покупатель --- человек в возрасте от 18 до 45 лет, активно занимающийся
спортом или фитнесом, имеющий опыт онлайн-покупок и ожидающий интуитивно
понятного интерфейса, быстрого поиска нужного товара и удобного оформления
заказа. Ключевыми требованиями для данной роли являются: наглядный каталог с
фильтрацией по категориям и характеристикам, корзина с возможностью изменения
состава, а также личный кабинет с историей заказов и редактированием профиля.

Администратор, как правило, является сотрудником компании со средним уровнем
технической грамотности. Его основные задачи --- своевременное обновление
ассортимента, контроль поступающих заказов и изменение их статусов по мере
выполнения. Интерфейс административной панели должен обеспечивать быстрый
доступ к ключевым операциям и минимизировать вероятность ошибок при управлении
данными.

Для каждой роли проведено профилирование по методу персонажей (Personas) для
минимизации когнитивной нагрузки на интерфейс.


\subsection{Описание функциональности приложения}

Функциональные возможности системы определяются через взаимодействие пользователя
с системой. Основной бизнес-кейс представлен в виде \textit{User Story}: «Я, как
покупатель, хочу добавить товары в корзину и оформить заказ, чтобы получить
выбранные спортивные товары с доставкой на указанный адрес».

На рисунке~\ref{fig:use_case} представлена диаграмма вариантов использования
(Use Case).

% Раскомментируй когда будет готов файл use_case_diagram.png
\begin{figure}[ht!]
    \centering
    \includegraphics[width=0.8\textwidth]{use_case_diagram}
    \caption{Диаграмма вариантов использования системы}
    \label{fig:use_case}
\end{figure}


\subsection{Проектирование модели данных}

При выборе архитектуры базы данных был проведён сравнительный анализ эффективности
различных решений, опирающийся на данные о производительности современных
СУБД~\cite{petrov2024}.

Структура таблицы заказов представлена в табл.~\ref{tab:db_orders}.

\begin{longtable}{|l|l|p{3.5cm}|p{5.5cm}|}
    \caption{Структура таблицы заказов (\texttt{orders})} \label{tab:db_orders} \\

    \hline
    \textbf{Поле} & \textbf{Тип} & \textbf{Ограничение} & \textbf{Описание} \\ \hline
    \endfirsthead

    \tabcont{\ref{tab:db_orders}}
    \textbf{Поле} & \textbf{Тип} & \textbf{Ограничение} & \textbf{Описание} \\ \hline
    \endhead

    \hline
    \endfoot

    \hline
    \endlastfoot

    id             & UUID      & PRIMARY KEY   & Уникальный идентификатор заказа        \\ \hline
    user\_id       & UUID      & FOREIGN KEY   & Ссылка на таблицу пользователей        \\ \hline
    status         & VARCHAR   & NOT NULL      & Статус (новый, в обработке, доставлен) \\ \hline
    total\_price   & NUMERIC   & NOT NULL      & Итоговая стоимость заказа, руб.        \\ \hline
    address        & TEXT      & NOT NULL      & Адрес доставки                         \\ \hline
    contact\_phone & VARCHAR   & NOT NULL      & Телефон получателя                     \\ \hline
    comment        & TEXT      & NULL          & Комментарий к заказу                   \\ \hline
    created\_at    & TIMESTAMP & DEFAULT NOW() & Время оформления заказа                \\ \hline
    updated\_at    & TIMESTAMP & DEFAULT NOW() & Время последнего изменения             \\ \hline
\end{longtable}

Итоговая стоимость заказа рассчитывается следующим образом:

\begin{equation}
    C_{total} = \sum_{i=1}^{n} (Q_i \cdot P_i) \cdot (1 - D) + S_{delivery},
    \label{eq:total_cost}
\end{equation}

\noindent где $Q_i$ --- количество единиц $i$-го товара в заказе;
$P_i$ --- цена единицы $i$-го товара, руб.;
$D$ --- коэффициент скидки (от 0 до 1);
$S_{delivery}$ --- стоимость доставки, руб.

На рисунке~\ref{fig:er_diagram} представлена ER-диаграмма базы данных,
отражающая сущности системы и связи между ними.

% Раскомментируй когда будет готов файл er_diagram.png
\begin{figure}[ht!]
    \centering
    \includegraphics[width=0.95\textwidth]{er_diagram}
    \caption{ER-диаграмма базы данных интернет-магазина}
    \label{fig:er_diagram}
\end{figure}

Логическая модель данных включает следующие основные сущности:
\begin{itemize}
    \item \textbf{users} --- пользователи системы (покупатели и администраторы);
    \item \textbf{products} --- товары каталога с привязкой к категориям;
    \item \textbf{categories} --- иерархический справочник категорий товаров;
    \item \textbf{orders} --- заказы, оформленные пользователями;
    \item \textbf{order\_items} --- позиции заказа (связующая таблица
          между заказами и товарами);
    \item \textbf{cart\_items} --- товары, добавленные пользователем в корзину.
\end{itemize}

Физическая схема базы данных, представленная на рисунке~\ref{fig:physical_schema},
специфицирует конкретные типы данных, индексы и внешние ключи применительно
к выбранной СУБД PostgreSQL.

% Раскомментируй когда будет готов файл physical_schema.png
\begin{figure}[ht!]
    \centering
    \includegraphics[width=0.95\textwidth]{physical_schema}
    \caption{Физическая схема базы данных}
    \label{fig:physical_schema}
\end{figure}


\subsection{Разработка макетов страниц}

В рамках проектирования пользовательского интерфейса были разработаны
низкодетализированные макеты (wireframes) основных страниц приложения.
Макеты отражают структуру интерфейса, расположение ключевых элементов
и пользовательский путь от выбора товара до оформления заказа.

На рисунке~\ref{fig:wf1} представлен макет главной страницы интернет-магазина,
включающий навигационную панель, hero-баннер, блок категорий и карточки
популярных товаров.

\begin{figure}[ht!]
    \centering
    \includegraphics[width=0.9\textwidth]{images/wireframe_1}
    \caption{Макет главной страницы}
    \label{fig:wf1}
\end{figure}

На рисунке~\ref{fig:wf2} показана страница каталога товаров с боковой
панелью фильтрации по цене, бренду и размеру, а также сеткой карточек
с пагинацией.

\begin{figure}[ht!]
    \centering
    \includegraphics[width=0.9\textwidth]{images/wireframe_2}
    \caption{Макет страницы каталога товаров}
    \label{fig:wf2}
\end{figure}

Карточка товара (рисунок~\ref{fig:wf3}) содержит галерею изображений,
элементы выбора размера и цвета, управление количеством и кнопку
добавления в корзину.

\begin{figure}[ht!]
    \centering
    \includegraphics[width=0.9\textwidth]{images/wireframe_3}
    \caption{Макет карточки товара}
    \label{fig:wf3}
\end{figure}

На рисунке~\ref{fig:wf4} представлен макет корзины с перечнем выбранных
товаров, управлением количеством и блоком сводки стоимости заказа.

\begin{figure}[ht!]
    \centering
    \includegraphics[width=0.9\textwidth]{images/wireframe_4}
    \caption{Макет страницы корзины}
    \label{fig:wf4}
\end{figure}

Страница оформления заказа (рисунок~\ref{fig:wf5}) реализована в виде
многошагового мастера с прогресс-баром, формой ввода данных получателя,
адресом доставки и итоговой сводкой.

\begin{figure}[ht!]
    \centering
    \includegraphics[width=0.9\textwidth]{images/wireframe_5}
    \caption{Макет страницы оформления заказа}
    \label{fig:wf5}
\end{figure}

После успешного оформления пользователь переходит на страницу
подтверждения (рисунок~\ref{fig:wf6}), где отображается номер заказа,
его статус и кнопки для дальнейших действий.

\begin{figure}[ht!]
    \centering
    \includegraphics[width=0.9\textwidth]{images/wireframe_6}
    \caption{Макет страницы подтверждения заказа}
    \label{fig:wf6}
\end{figure}

Личный кабинет пользователя (рисунок~\ref{fig:wf7}) включает боковое
меню навигации, аватар с контактными данными и список оформленных заказов
с указанием их статусов.

\begin{figure}[ht!]
    \centering
    \includegraphics[width=0.9\textwidth]{images/wireframe_7}
    \caption{Макет личного кабинета пользователя}
    \label{fig:wf7}
\end{figure}

Административная панель (рисунок~\ref{fig:wf8}) предоставляет сотруднику
магазина доступ к ключевым показателям (количество заказов, выручка,
число товаров и пользователей), а также таблицу последних заказов
с возможностью управления их статусами.

\begin{figure}[ht!]
    \centering
    \includegraphics[width=0.9\textwidth]{images/wireframe_8}
    \caption{Макет административной панели}
    \label{fig:wf8}
\end{figure}